%! Compressing Images by the SVD — Unit 7, Lesson 7.2 — Guided Notes
\documentclass[10pt]{article}
\usepackage{linalg-article}
\usepackage{linalg-boxes}
\usetikzlibrary{arrows.meta}
\newcommand{\TallMath}[1]{$\displaystyle #1\rule[-1.4em]{0pt}{3.2em}$}
\newcommand{\uu}{\mathbf{u}}
\newcommand{\vv}{\mathbf{v}}
\newcommand{\ee}{\mathbf{e}}
\newcommand{\zero}{\mathbf{0}}
\newcommand{\T}{^{\top}}

\begin{document}

\pageheader{Unit 7, Lesson 7.2}{Guided Notes}
\namedateperiod

\vspace{0.12in}

\begin{objectivebox}
\textbf{By the end of this lesson, I will be able to\dots}
\begin{itemize}\setlength{\itemsep}{2pt}
  \item build a rank-1 matrix as an outer product $\uu\vv\T$ (a column times a row);
  \item write a matrix as a sum of rank-1 layers $A=\sigma_1\uu_1\vv_1\T+\sigma_2\uu_2\vv_2\T+\cdots$ and keep only the biggest;
  \item explain why keeping the largest singular values gives the best low-storage approximation of an image.
\end{itemize}
\end{objectivebox}

\vspace{0.06in}

\begin{vocabbox}
\small
Fill in each term as we build it together.
\vspace{2pt}
\termblanklong{Outer product $\uu\vv\T$ (a column times a row --- a rank-1 matrix)}
\termblanklong{The SVD as layers: $A=\sigma_1\uu_1\vv_1\T+\sigma_2\uu_2\vv_2\T+\cdots$}
\termblanklong{Rank-$k$ approximation $A_k$ (keep the $k$ biggest layers)}
\end{vocabbox}

\begin{hookbox}
A grayscale photo is just a big matrix of pixel-brightness numbers --- a $1000\times1000$ photo holds a
million of them. Storing every number is expensive.
\begin{itemize}\setlength{\itemsep}{1pt}
  \item The SVD splits a matrix into ordered layers $\sigma_1\uu_1\vv_1\T,\ \sigma_2\uu_2\vv_2\T,\dots$ Which layers do you think matter most --- and why? \writeline
  \item If you could keep only a \emph{few} layers and throw the rest away, which would you keep to lose the least? \writeline
\end{itemize}
\end{hookbox}

\begin{notesbox}{1. The Outer Product: a Column Times a Row}
The SVD is built from the ingredients $\sigma,\uu,\vv$ of Lesson~7.1. To use them we need one new move:
multiply a \emph{column} $\uu$ by a \emph{row} $\vv\T$. The result is a whole matrix, called an
\textbf{outer product}:
\[
  \uu\vv\T=\begin{bmatrix} u_1\\ u_2 \end{bmatrix}\begin{bmatrix} v_1 & v_2 \end{bmatrix}
  =\begin{bmatrix} u_1v_1 & u_1v_2\\ u_2v_1 & u_2v_2 \end{bmatrix}.
\]
Try it with $\uu=\begin{bsmallmatrix} 1\\2 \end{bsmallmatrix}$ and $\vv\T=\begin{bsmallmatrix} 3 & 1 \end{bsmallmatrix}$:
\[
  \uu\vv\T=\begin{bmatrix}\rule{0pt}{1.1em}\blank{0.9cm} & \blank{0.9cm}\\[2pt]\blank{0.9cm} & \blank{0.9cm}\end{bmatrix}.
\]
Notice: \emph{every column} is a multiple of $\uu$ (and every row a multiple of $\vv\T$). A matrix this
simple --- made from one column and one row --- is called a \textbf{rank-1} matrix. It stores in only
$\blank{1.2cm}+\blank{1.2cm}$ numbers (the entries of $\uu$ and $\vv$), not the full $2\times2=4$.
\end{notesbox}

\begin{notesbox}{2. The SVD is a Sum of Rank-1 Layers}
From Lesson~7.1, $A=U\Sigma V\T$. Multiplying that out groups $A$ into one rank-1 outer product for each
singular value:
\[
  A=\sigma_1\,\uu_1\vv_1\T+\sigma_2\,\uu_2\vv_2\T+\cdots+\sigma_r\,\uu_r\vv_r\T .
\]
Each \textbf{layer} $\sigma_i\uu_i\vv_i\T$ is a rank-1 matrix scaled by its singular value, and the layers
are ordered from the biggest $\sigma$ to the smallest.

\vspace{2pt}
\textbf{Worked on today's matrix} $A=\begin{bsmallmatrix} 5 & 3\\ 3 & 5 \end{bsmallmatrix}$. From the
Warm-Up: $\sigma_1=8,\ \sigma_2=2$, with $\uu_1=\vv_1=\tfrac1{\sqrt2}\begin{bsmallmatrix} 1\\1 \end{bsmallmatrix}$
and $\uu_2=\vv_2=\tfrac1{\sqrt2}\begin{bsmallmatrix} 1\\-1 \end{bsmallmatrix}$ (this $A$ is symmetric, so
$\uu_i=\vv_i$). Build the two layers --- the two $\tfrac1{\sqrt2}$'s multiply to $\tfrac12$:
\[
  \sigma_1\uu_1\vv_1\T=8\cdot\tfrac12\begin{bmatrix} 1 & 1\\ 1 & 1 \end{bmatrix}
  =\begin{bmatrix}\rule{0pt}{1.1em}\blank{0.7cm} & \blank{0.7cm}\\[2pt]\blank{0.7cm} & \blank{0.7cm}\end{bmatrix},\qquad
  \sigma_2\uu_2\vv_2\T=2\cdot\tfrac12\begin{bmatrix} 1 & -1\\ -1 & 1 \end{bmatrix}
  =\begin{bmatrix}\rule{0pt}{1.1em}\blank{0.7cm} & \blank{0.7cm}\\[2pt]\blank{0.7cm} & \blank{0.7cm}\end{bmatrix}.
\]
Add the layers and you get $A$ back:
$\begin{bsmallmatrix} 4 & 4\\ 4 & 4 \end{bsmallmatrix}+\begin{bsmallmatrix} 1 & -1\\ -1 & 1 \end{bsmallmatrix}
=\begin{bsmallmatrix}\rule{0pt}{1em}\blank{0.6cm} & \blank{0.6cm}\\[1pt]\blank{0.6cm} & \blank{0.6cm}\end{bsmallmatrix}=A.$
\end{notesbox}

\begin{notesbox}{3. Compression: Keep the Biggest Layers}
Because the layers are ordered $\sigma_1\ge\sigma_2\ge\cdots$, the \emph{first} layer carries the most and
the last carries the least. Throwing away the small-$\sigma$ layers barely changes $A$. Keeping just the
first $k$ layers gives the \textbf{rank-$k$ approximation}:
\[
  A_k=\sigma_1\uu_1\vv_1\T+\cdots+\sigma_k\uu_k\vv_k\T .
\]
\textbf{Eckart--Young fact:} $A_k$ is the \emph{closest} rank-$k$ matrix to $A$ --- no simpler matrix
approximates $A$ better. So the best way to store $A$ cheaply is to keep its biggest singular values.

\vspace{2pt}
For our matrix, the rank-1 approximation is just layer~1:
$A_1=\begin{bsmallmatrix}\rule{0pt}{1em}\blank{0.6cm} & \blank{0.6cm}\\[1pt]\blank{0.6cm} & \blank{0.6cm}\end{bsmallmatrix}$
--- it differs from $A$ by only the small dropped layer $\begin{bsmallmatrix} 1 & -1\\ -1 & 1 \end{bsmallmatrix}$.

\begin{center}
\begin{tikzpicture}[scale=1, font=\footnotesize]
  \def\bw{0.62}
  % bar centers and heights (a decaying singular-value spectrum)
  \foreach \x/\h/\c/\lab in {0.7/2.9/burgundy/$\sigma_1$, 1.6/1.85/burgundy/$\sigma_2$,
        2.5/1.05/blushmid/$\sigma_3$, 3.4/0.62/blushmid/$\sigma_4$, 4.3/0.38/blushmid/$\sigma_5$,
        5.2/0.22/blushmid/$\sigma_6$} {
    \fill[\c] (\x-\bw/2,0) rectangle (\x+\bw/2,\h);
    \draw[charcoal] (\x-\bw/2,0) rectangle (\x+\bw/2,\h);
    \node[below] at (\x,-0.02) {\lab};
  }
  \draw[-{Latex[length=2mm]}, thick] (0.1,0)--(5.9,0);
  % keep/drop divider
  \draw[dashed, burgundy, thick] (2.05,-0.35) -- (2.05,3.2);
  \node[burgundy] at (1.15,3.05) {\scriptsize\textbf{keep}};
  \node[charcoal] at (3.8,3.05) {\scriptsize drop};
\end{tikzpicture}\\[-1pt]
{\scriptsize A real image has many singular values that fall off fast. Keep the few big ones (burgundy),
drop the small tail --- a close copy in a fraction of the space.}
\end{center}

\vspace{2pt}
\textbf{How much did we keep?} The total ``size'' of $A$ (the sum of the squares of all its entries) splits
across the singular values: $\sigma_1^2+\sigma_2^2=64+4=68$ (Warm-Up~3). Layer~1 alone holds $\sigma_1^2=64$,
which is $\dfrac{64}{68}\approx\blank{1.2cm}\%$ of the total --- \emph{one} rank-1 layer already captures
about $94\%$ of the matrix.
\end{notesbox}

\begin{notesbox}{4. Why It Saves So Much --- and the Road Ahead}
\textbf{The storage win (a real image).} A $1000\times1000$ grayscale photo is $1{,}000{,}000$ numbers. Each
rank-1 layer needs only $\sigma$ plus $\uu$ plus $\vv$, i.e. $1+1000+1000=2001$ numbers. Keeping the
$20$ biggest layers stores about $20\times2001\approx\blank{2.2cm}$ numbers --- roughly
\blank{1.4cm}\% of the original --- and for real photos the singular values drop off so fast that the
picture looks nearly identical.

\vspace{2pt}
\textbf{Why the biggest $\sigma$'s?} Ordering by $\sigma$ means the first layers carry the most ``energy''
($\sigma^2$), so dropping the tail loses the least. And \emph{every} matrix has this layered SVD
(Lesson~7.1), so \emph{every} image compresses this way.

\vspace{2pt}
\textbf{The road ahead.} \textbf{7.3} principal component analysis --- the same ``keep the biggest
$\sigma$'' idea applied to \emph{data} instead of pixels\; $\bullet$\; \textbf{7.4} the victory of orthogonality.
\end{notesbox}

\begin{practicebox}
Show your work.
\begin{enumerate}[leftmargin=1.6em, itemsep=4pt]
  \item Compute the outer product $\begin{bsmallmatrix} 1\\3 \end{bsmallmatrix}\begin{bsmallmatrix} 2 & 2 \end{bsmallmatrix}$, and explain why it is a rank-1 matrix. \writeline
  \item The notes' rank-1 approximation is $A_1=\begin{bsmallmatrix} 4 & 4\\ 4 & 4 \end{bsmallmatrix}$. Compare it to $A=\begin{bsmallmatrix} 5 & 3\\ 3 & 5 \end{bsmallmatrix}$: by how much is each entry off, and what fraction of the ``energy'' $\sigma_1^2+\sigma_2^2$ does $A_1$ keep? \writeline
  \item In one sentence: why does keeping the \emph{largest} singular values give the best low-storage copy of an image? \writeline
\end{enumerate}
\end{practicebox}

\end{document}
